\documentclass[a4paper]{article}
\usepackage[pdftex]{graphicx}
\usepackage[utf8]{inputenc}
\usepackage{enumerate}
\usepackage{icomma}
\usepackage{amssymb}
\usepackage{tikz}
\usepackage{href-ul}
\hypersetup{
	colorlinks=true,
	linkcolor=blue,
	urlcolor=blue}
\usepackage{geometry}
\geometry{a4paper, top=15mm, left=15mm, right=15mm, bottom=15mm,
	headsep=10mm, footskip=12mm}

\begin{document}
\begin{tikzpicture}(10,3)
	\draw[ultra thick](2,0) --(10,0) -- (10,1.3) --(2,1.3) -- (2,0);
	\draw[fill=black](2,0)-- (0,.6) -- (2,.6) -- (2,0);
	\node at (6,.6) {\large Lösungsblatt von \href{https://www.okuyakl.de}{www.okuyakl.de}};
\end{tikzpicture}
\vspace{0.5 cm}

\noindent
\fbox{
	\begin{minipage}{0.48\textwidth}
	\noindent{\bf Aufgabe 1.}\\
	\begin{itemize}
		\item Zähler enthält Term $(x-2)$
		\item Nenner enthält Term $(x-3)^2$
		\item Zählergrad = Nennergrad und die Faktoren vor den höchsten Potenzen müssen 3 ergeben:
	\end{itemize}
	$$\Rightarrow \quad f(x)= {3x(x-2)\over (x-3)^2}$$	
	\end{minipage}
}
\hfill
\fbox{
	\begin{minipage}{0.48\textwidth}
	\noindent{\bf Aufgabe 2.}\\
	$$
	\renewcommand{\arraystretch}{2}
	\begin{array}{rcccccl}
	{f(x) -f(2) \over x-2} &=& {{1\over x}+3 - 3,5 \over x-2} &=& {{1\over x }-{1\over 2} \over x-2}\\
	&=&  {{2-x \over 2x} \over x-2} &=& -{1 \over 2x} \\
	f'(2)&=& \lim\limits_{x \to 2}  -{1 \over 2x} &=& -{1 \over 4}
	\end{array}
	$$	
	\end{minipage}
}

\noindent
\fbox{
	\begin{minipage}{0.48\textwidth}
	\noindent{\bf Aufgabe 3. a)}\\
	Wir setzen den Nenner gleich null:
	$$
	\renewcommand{\arraystretch}{2}
	\begin{array}{rcll}
	x^2-4 &=&0 \\
	(x-2)(x+2) &=& 0 
	\end{array}
	$$
	Wir erkennen in der faktorisierten Form zwei einfache Definitionslücken bei $2$ und $-2$. Der Zähler hat an diesen Stellen keine Nullstelle, also sind dies zwei Polstellen erster Ordnung (ohne VZW).
	$$\mathbb{D}_f=\mathbb{R} \setminus \{-2;2\}$$
	
	\noindent{\bf Aufgabe 3. b)}\\
	$$f(-x)= {(-x)^3 \over (-x)^2-4}= -{x^3 \over x^2-4}=-f(x) $$
	Die Funktion ist punktsymmetrisch.
	
	\noindent{\bf Aufgabe 3. c)}\\
	Senkrechte Asymptoten:
	$$x=-2 \quad \textnormal{und} \quad x=2$$
	Zählergrad $>$ Nennergrad; Wir führen die Polynomdivision durch:
	$$x^3:(x^2-4)= x +{4x \over x^2-4}$$
	Daraus folgt, dass 
	$$y=x$$
	schräge Asymptote ist.	
	\end{minipage}
}
\hfill
\fbox{
	\begin{minipage}{0.48\textwidth}
	\noindent{\bf Aufgabe 3. d)}\\
	Wir bilden die erste Ableitung:
	$$f'(x)={(x^2-4)\cdot 3x^2-x^3\cdot2x \over (x^2-4)^2} = {x^4-12x^2 \over  (x^2-4)^2}$$
	
	\noindent{\bf Aufgabe 3. e)}\\
	Punkt P:
	$$f(4)={16\over 3} \quad \Rightarrow P(4|{16\over 3})$$
	Steigung m:
	$$m=f'(4)={4\over 9} \quad \Rightarrow y=mx+t$$
	y-Achsenabschnitt t:
	$$\Rightarrow {16\over 3} = {4\over 9}\cdot 4 +t  \quad \Rightarrow t= {32 \over 9}$$
	Funktionsterm Tangente:
	$$y={4\over 9}x + {32 \over 9}$$
	Schnittwinkel x-Achse:
	$$m=\tan(\alpha) \quad \Rightarrow \alpha=24,0^\circ$$
		
	\end{minipage}
}


\noindent
\fbox{
	\begin{minipage}{0.48\textwidth}
	\noindent{\bf Aufgabe 4. a) b)}\\
	\begin{itemize}
		\item Doppelte Nullstelle:$x=0 \quad \Rightarrow x^2$ im Zähler
		\item Definitionslücken:$x=\pm 2 \quad \Rightarrow (x^2-4)$ im Nenner
		\item Weitere Asymptote: $y=2 \quad \Rightarrow $ Zählergrad =Nennergrad; Faktor 2.
	\end{itemize}
	$$f(x)={2x^2 \over x^2-4}$$	
	\end{minipage}
}
\hfill
\fbox{
	\begin{minipage}{0.48\textwidth}
	
	\noindent{\bf Aufgabe 4. c)}\\
	$$f(-x)={2(-x)^2 \over (-x)^2-4}={2x^2 \over x^2-4}=f(x)$$
		
	\end{minipage}
}



\noindent
\begin{minipage}{0.6\textwidth}
	\noindent{\bf Aufgabe 5. a)}\\
	Faktorisieren des Nenners liefert:
	$$4x^2(x^2+x-6)=4x^2(x+3)(x-2)$$
	Die Definitionsmenge ist damit:
	$$\mathbb{D}_f=\mathbb{R} \setminus \{-3; 0; 2\}$$
	Nullstellen sind Nullstellen des Zählers:
	$$(x^2-9)=(x+3)(x-3)=0 \quad \Rightarrow x_1=3 \quad x_2=-3$$
\end{minipage}
\hfill
\begin{minipage}{0.4\textwidth}
	\includegraphics[width=6 cm]{fumf061}
\end{minipage}
\vspace{0.5 cm}

\noindent
Der Linearfaktor $(x+3)$ kommt sowohl im Zähler als auch im Nenner vor.
Wir kürzen ihn und führen die Grenzwertbetrachtung durch:
$$\lim\limits_{x\to -3} {x-3 \over 4x^2(x-2)}={-6 \over -180}={1\over 30}$$
Dieser Grenzwert ist beidseitig, hierbei handelt es sich um eine stetig hebbare Definitionslücke.
\vspace{0.5 cm}

\noindent
\begin{minipage}{0.6\textwidth}
	\noindent{\bf Aufgabe 6. a)}\\
	Wir bringen alles auf einen gemeinsamen Nenner:
	$${x^2-2x+1 \over x^2}=0 \quad \Rightarrow x^2-2x+1 = (x-1)^2=0$$
	Hieraus folgt, dass $x=1$ einzige, doppelte Nullstelle ist.
	$$\lim\limits_{x\to0+}{x^2-2x+1 \over x^2}="{1 \over +0}" = + \infty$$
	$$\lim\limits_{x\to0-}{x^2-2x+1 \over x^2}="{1 \over +0}" = + \infty$$
	
\end{minipage}
\hfill
\begin{minipage}{0.4\textwidth}
	\noindent{\bf Aufgabe 6. b)}\\
	\includegraphics[width=7 cm]{seks061}
\end{minipage}

\noindent
\begin{minipage}{0.6\textwidth}
	\noindent{\bf Aufgabe 7. a)}\\
	Nullstellen des Zählers = Nullstellen
	$$x^2+3x=0 \quad \Rightarrow x_1=0 \quad x_2=-3$$
	Zwei einfache Nullstellen.\\ 
	Nullstellen des Nenners = Definitionslücken
	$$4(x-1)=0 \quad \Rightarrow x=1$$ 
	$$\mathbb{D}=\mathbb{R} \setminus \{1\}$$
	$$\lim\limits_{x\to1-}{x^2+3x \over 4(x-1)}="{4 \over -0}" = - \infty$$
	$$\lim\limits_{x\to1+}{x^2+3x \over 4(x-1)}="{4 \over +0}" = + \infty$$
\end{minipage}
\hfill
\begin{minipage}{0.4\textwidth}
	\noindent{\bf Aufgabe 7. b)}\\
	\includegraphics[width=7 cm]{siebn061}
\end{minipage}
\vspace{0.5 cm}

Die schräge Asymptote finden wir mit Polynomdivision:
$$(x^2+3x):(4x-4)={1\over 4}x+1+ {1 \over x-1} \quad \Rightarrow y={1\over 4}x+1$$

\begin{center}
	\includegraphics[width=7 cm]{../../viecher/endcomic.pdf}
	
	Hier geht es zurück zum \href{https://www.okuyakl.de/math/m11gbrL061/aa061.pdf}{Aufgabenblatt}
\end{center}

\end{document}

